
\documentclass[a4paper,10pt]{report}

%\usepackage[landscape]{geometry}
\usepackage[cm-default]{fontspec}
\setromanfont{FreeSerif}
\setsansfont{FreeSans}
\setmonofont{FreeMono}
\usepackage[utf8x]{inputenc}
\usepackage{fontspec}
\usepackage{xunicode}
\usepackage{xltxtra}
\usepackage{xgreek}
\usepackage{amsmath}
\usepackage{unicode-math}
\usepackage{ulem}
\usepackage{color}
\usepackage{verbatim}
\usepackage{nopageno}
\usepackage{graphicx}
\usepackage{textpos}
\setlength{\TPHorizModule}{1cm}
\setlength{\TPVertModule}{1cm}
%\usepackage[colorgrid,texcoord]{eso-pic}
\usepackage[outline]{contour}
\usepackage{wrapfig}
\usepackage{url}
\usepackage{color}
\usepackage{tikz}
\usepackage{fancybox,fancyhdr}
\usepackage{subfigure}
\usepackage{pstricks}
\usepackage{epsfig}
\usepackage{multicol}
\usepackage{listings}
\usepackage{enumerate}
\usepackage{hyperref}
\hypersetup{
  bookmarks=true,
  bookmarksopen=true,
  pdfborder=false,
  pdfpagemode=UseNone,
  raiselinks=true,
  pdfhighlight={/P},
  colorlinks,
  citecolor=black,
  filecolor=black,
  linkcolor=black,
  urlcolor=black
}


\usepackage{fontspec}
\usepackage{xunicode}
\usepackage{xltxtra}


% Margins
%\setlength{\textwidth}{24cm}

%\setlength{\voffset}{0in}
%\setlength{\textheight}{6.8in}
% Colors
\definecolor{cbrown}{rgb}{0.49,0.24,0.07}
\definecolor{cmpez}{rgb}{0.92,0.88,0.79}
\definecolor{cmpez2}{rgb}{0.62,0.33,0.34}
\definecolor{cwhite}{rgb}{1,1,1}
\definecolor{cblack}{rgb}{0,0,0}
\definecolor{cred}{rgb}{0.9,0.15,0.15}
\definecolor{clblue}{rgb}{0.57,0.85,0.97}
\definecolor{corange}{rgb}{0.86,0.49,0.18}
\definecolor{clorange}{rgb}{0.95,0.84,0.61}
\definecolor{cyellow}{rgb}{1,0.95,0.16}
\definecolor{cgreen}{rgb}{0.69,0.8,0.31}
\definecolor{clbrown}{rgb}{0.78,0.67,0.43}
\definecolor{cmagenta}{rgb}{0.79,0.05,0.54}
\definecolor{cgray}{rgb}{0.76,0.73,0.63}
\definecolor{bluesite}{rgb}{0.05,0.43,0.67}
\definecolor{redsite}{rgb}{0.90,0.15,0.15}
\definecolor{backsite}{rgb}{0.07,0.07,0.07}
% ----------- from Costas -----------------------------

%---------------   Main Fonts    -----------------%   

% \setmainfont[Mapping=tex-text]{Calibri}
% \setmainfont[Mapping=tex-text]{Times New Roman}
%\setmainfont[Mapping=tex-text]{Droid Serif}
%\setmainfont[Mapping=tex-text]{Cambria}
%\setmainfont[Mapping=tex-text]{cm-unicode}
% \setmainfont[Mapping=tex-text]{Gentium}
% \setmainfont[Mapping=tex-text]{GFS Didot}
% \setmainfont[Mapping=tex-text]{Comic Sans MS}
 \setmainfont[Mapping=tex-text]{Ubuntu}
% \setmainfont[Mapping=tex-text]{Myriad Pro}
% \setmainfont[Mapping=tex-text]{CMU Concrete}
%\setmainfont[Mapping=tex-text]{DejaVu Sans}
%\setmainfont[Mapping=tex-text]{KerkisSans}
%\setmainfont[Mapping=tex-text]{KerkisCaligraphic}
% \setmainfont[Mapping=tex-text]{Segoe Print}
% \setmainfont[Mapping=tex-text]{Gabriola}

%------------------------------------------------%

% ------------ Mathematics Fonts ----------------%
% \setmathfont{Asana-Math.ttf}
% -----------------------------------------------%
% Misc
\author{Κασωτάκη Ε. - Σμαραγδάκης Κ.}
\title{www.math24.gr}
% ---------------- Formation --------------------%


\setlength\topmargin{-0.7cm}
\addtolength\topmargin{-\headheight}
\addtolength\topmargin{-\headsep}
\setlength\textheight{26cm}
\setlength\oddsidemargin{-0.54cm}
\setlength\evensidemargin{-0.54cm}
%\setlength\marginparwidth{1.5in}
\setlength\textwidth{17cm}

\RequirePackage[avantgarde]{quotchap}
\renewcommand\chapterheadstartvskip{\vspace*{0\baselineskip}}
\RequirePackage[calcwidth]{titlesec}
\titleformat{\section}[hang]{\bfseries}
{\Large\thesection}{12pt}{\Large}[{\titlerule[0.9pt]}]
%--------------------------------------------------%
%\usepackage{draftwatermark}
%\SetWatermarkText{www.math24.gr}
%\SetWatermarkLightness{0.9}
%\SetWatermarkScale{0.6}
%--------------------------------------------------%


\usepackage{xcolor}
\usepackage{amsthm}
\usepackage{framed}
\usepackage{parskip}

\colorlet{shadecolor}{bluesite!20}

\newtheorem*{orismos}{Ορισμός}

\newenvironment{notation}
  {\begin{shaded}\begin{theorem}}
  {\end{theorem}\end{shaded}}




% Document begins
\begin{document}
\pagestyle{fancy}
\fancyhead{}
\fancyfoot{}
\renewcommand{\headrulewidth}{0pt}
\renewcommand{\footrulewidth}{0pt}


\fancyhead[LO,LE]{
 %\textblockcolor{backsite}
 %\begin{textblock}{5}(-2,-0.55)
  %\rule{0cm}{1cm}
 %\end{textblock}
 \textblockcolor{bluesite}
 \begin{textblock}{5}(-1.5,-0.55)
  \rule{0cm}{1cm}
 \end{textblock}
 %\textblockcolor{bluesite}
 %\begin{textblock}{14}(5.5,-0.55)
 % \rule{0cm}{1cm}
 %\end{textblock}
 \begin{textblock}{0}(-1,-0.25)
 \color{cwhite} \begin{Large}www.math24.gr\end{Large}
 \end{textblock}
\begin{textblock}{0}(16,-0.4)
 \color{cwhite} \includegraphics[height=1.5cm]{math24_logo.png}
 \end{textblock}
\textblockcolor{white}
\begin{textblock}{17}(-1,28)
 \color{backsite} \begin{small}Copyright \textcopyright 2012, Κασωτάκη Ε.(ikasotaki@gmail.com) - Σμαραγδάκης Κ.(kesmarag@gmail.com)\end{small}
 \end{textblock}
}
 
\begin{shaded}
\begin{center}
\huge \textbf{Φυλλάδιο Ασκήσεων}\\
%Πρόσθεση ρητών αριθμών
\end{center} 
\textbf{Μαθηματικά Α' Γυμνασίου} \hfill \textbf{Ημερομηνία Παράδοσης : \hspace{2em} }
\subsection*{Ονοματεπώνυμο :\hfill  \hspace{5em}}
\end{shaded}
\vspace{2em}
\begin{itemize}
 \item Πολλαπλασιασμός φυσικών αριθμών
 \item Πολλαπλασιασμός με $10,100,\cdots$
 \item Ιδιότητες του πολλαπλασιασμού
 \item Επιμεριστική ιδιότητα
\end{itemize}


\section*{Θεωρία - Πολλαπλασιασμός Φυσικών Αριθμών \hfill \small{}}
\textbf{Πολλαπλασιασμός} είναι η πράξη με την οποία από δύο φυσικούς αριθμούς $α$ και $β$ (τους παράγοντες), 
βρίσκουμε ένα τρίτο φυσικό αριθμό $γ$ που είναι το γινόμενό τους. Δηλαδή $α\cdot β=γ$\\
\textbf{π.χ} $3\cdot 5 =15$ (παράγοντες: $3$ και $5$, γινόμενο: $15$)

\section*{Άσκηση 1  \hfill \small{10 μονάδες}}
Να υπολογίσετε τα παρακάτω γινόμενα: 
\begin{enumerate} [1)]
 \item $33\cdot 11$
 \item $2\cdot 18$
 \item $6\cdot 51$
 \item $13\cdot 4$
 \item $12\cdot 12$
\end{enumerate}

\section*{Θεωρία - Πολλαπλασιασμός με $10,100,\cdots$ \hfill \small{}}
Για να πολλαπλασιάσουμε ένα αριθμό επί $10,100,1000,\cdots$ γράφουμε στο τέλος του αριθμού τόσα μηδενικά όσα έχει 
κάθε φορά ο παράγοντας $10,100,\cdots$ \\
\textbf{π.χ} $71\cdot10=710$\\
\textbf{π.χ} $33\cdot100=3300$\\
\textbf{π.χ} $6\cdot1000=6000$\\
\textbf{π.χ} $768\cdot10000=7680000$

\section*{Άσκηση 2  \hfill \small{10 μονάδες}}
Να υπολογίσετε τα παρακάτω γινόμενα: 
\begin{enumerate} [1)]
 \item $3\cdot 1000$
 \item $49\cdot 10$
 \item $5\cdot 10000$
 \item $67\cdot 100$
 \item $3\cdot 100$
\end{enumerate}

\section*{Θεωρία - Ιδιότητες του πολλαπλασιασμού \hfill \small{}}
\textbf{Ιδιότητες του πολλαπλασιασμού:}
\begin{itemize}
 \item \textbf{Πολλαπλασιασμός επί $1$}: το $1$ όταν πολλαπλασιαστεί με ένα φυσικό αριθμό δεν τον μεταβάλλει,\\
        δηλαδή $α\cdot 1=α$ και $1\cdot α=α$\\
        \textbf{π.χ} $8\cdot1=8$\\
        \textbf{π.χ} $1\cdot16=16$
 \item \textbf{Αντιμεταθετική ιδιότητα: } Μπορούμε να αλλάξουμε τη σειρά των παραγόντων ενός γινομένου,\\
        δηλαδή $α\cdot β=β\cdot α$\\
       \textbf{π.χ} $7\cdot 11=11\cdot 7$ (γιατί $7\cdot11=77$ και $11\cdot 7=77$)
 \item \textbf{Προσεταιριστική ιδιότητα:} Μπορούμε να αντικαθιστούμε παράγοντες με το γινόμενό τους ή να αναλύουμε 
       ένα παράγοντα σε γινόμενο,\\
       δηλαδή $α\cdot (β\cdot γ)=(α\cdot β) \cdot γ)$\\
       \textbf{π.χ} $3\cdot (4\cdot 10)=(3\cdot 4)\cdot 10$ (γιατί $3\cdot(4\cdot10)=3\cdot40=120$ και 
                                                            $(3\cdot4)\cdot10=12\cdot10=120)$
\end{itemize}


\section*{Άσκηση 3  \hfill \small{}20  μονάδες}
Να αντιστοιχίσετε κάθε στοιχείο της αριστερής στήλης με ένα στοιχείο της δεξιάς στήλης
\begin{itemize}
\begin{multicols}{2}
 \item $671\cdot1$
 \item $336\cdot2$
 \item $2\cdot5\cdot11$
 \item $(3\cdot5)\cdot70$
 \item $10\cdot67$
 \item $3\cdot(5\cdot70)$
 \item $110$
 \item $(2\cdot5)\cdot67$
 \item $672$
 \item $671$
\end{multicols}
\end{itemize}


\section*{Θεωρία - Επιμεριστική Ιδιότητα  \hfill \small{}}
\begin{itemize}
 \item Επιμεριστική ιδιότητα του πολλαπλασιασμού ως προς την πρόσθεση: $α \cdot (β+γ)=α\cdot β+α\cdot γ$ \\
 \textbf{π.χ } $2\cdot(7+3)=2\cdot7+2\cdot3$ (γιατί $2\cdot (7+3)=2\cdot 10=20$ και $2\cdot7+2\cdot3=14+6=20$) \\ 
 \textbf{π.χ }Αντίστροφο της επιμεριστικής: $3\cdot1+3\cdot5=3\cdot(1+5)$  
(γιατί $3\cdot 1+3\cdot5=3+15=18$ και $3\cdot(1+5)=3\cdot6=18$) \\ 
 \item Επιμεριστική ιδιότητα του πολλαπλασιασμού ως προς την αφαίρεση: $α \cdot (β-γ)=α\cdot β-α\cdot γ$ \\
 \textbf{π.χ } $2\cdot(12-2)=2\cdot12-2\cdot2$ (γιατί $2\cdot (12-2)=2\cdot 10=20$ και $2\cdot12-2\cdot2=24-4=20$) \\ 
 \textbf{π.χ }Αντίστροφο της επιμεριστικής: $4\cdot3- 4\cdot2=4\cdot(3-2)$  
(γιατί $4\cdot 3 -4\cdot2=12-8=4$ και $4\cdot(3-2)=4\cdot1=4$) \\
\end{itemize}


\section*{Άσκηση 4  \hfill \small{20 μονάδες}}
Να υπολογίσετε τα παρακάτω γινόμενα χρησιμοποιώντας την επιμεριστική ιδιότητα: 
\begin{enumerate} [1)]
 \item $3\cdot (10-7)$
 \item $2\cdot (14-1)$
 \item $4\cdot (11+9)$
 \item $60\cdot (1+2)$
 \item $5\cdot (12-2)$
\end{enumerate}



\section*{Άσκηση 5  \hfill \small{20 μονάδες}}
Να υπολογίσετε τις παρακάτω πράξεις χρησιμοποιώντας το αντίστροφο της επιμεριστικής ιδιότητας: 
\begin{enumerate} [1)]
 \item $3\cdot4  +3\cdot 6$
 \item $2\cdot36  +2\cdot 64$
 \item $3\cdot21  +3\cdot 9$
 \item $71\cdot3  +71\cdot 7$
 \item $8\cdot6  +6\cdot 2$
\end{enumerate}



\section*{Άσκηση 6  \hfill \small{} 20 μονάδες}
Να αντιστοιχίσετε κάθε στοιχείο της αριστερής στήλης με ένα στοιχείο της δεξιάς στήλης
\begin{itemize}
\begin{multicols}{2}
 \item $78\cdot (2+1)$
 \item $91\cdot4+91\cdot 5$
 \item $13\cdot17-13\cdot7$
 \item $16\cdot(2+5)$
 \item $16\cdot(5-2)$
 \item $16\cdot 5-16\cdot2$
 \item $78\cdot 2+78\cdot1$
 \item $91\cdot(4+5)$
 \item $13\cdot(17-7)$
 \item $16\cdot2+16\cdot5$
\end{multicols}
\end{itemize}




 









%\let\thefootnote\relax\footnotetext{Επιμέλεια : Κασωτάκη Ειρήνη (ikasotaki@gmail.com),    Σμαραγδάκης Κώστας (kesmarag@gmail.com)} 


\end{document}
